% STATUS: rewritten 2026-10-08
% Sources read in full: Ritt 1950 (p. 1), Jacobson 1937 (pp. 206--207), Hochschild 1942 (pp. 677, 688--689), Kharchenko 1991 (pp. 6--8), Freudenburg 2017 (pp. 2, 13).
% Additions of 2026-11-07 moved from chapter 20: Dirac 1925 (pp. 648, 651) read in full.
\section{Definitions, first identities and the Leibniz formula}\label{sec:basic-theory:definitions}

We start from the product rule of calculus. The polynomial ring has the
derivative $d/dx$, and Section~\ref{sec:basic-theory:der-as-lie-algebra} gives its
derivations.
Matrix rings have no derivative that is obvious, and
Section~\ref{sec:basic-theory:inner-outer} shows where their derivations come
from. So the first task was to say which maps deserve the name, and to find out
what follows from the rule alone.

Ritt took the first step for fields \cite[p.~1]{Ritt1950}. He worked in
characteristic zero and asked for an operation $a \mapsto a'$ with
$(a+b)' = a'+b'$ and $(ab)' = ba' + ab'$. He saw at once that $0' = 0$ and
$1' = 0$. He named the elements with zero derivative: ``An element with zero
for derivative is called a constant'' \cite[p.~1]{Ritt1950}.

Jacobson took a wider view in 1937 \cite[p.~206]{Jacobson1937}. He wanted to study ``the set of operations
mapping an algebra on itself and having the formal character of derivation in
the field of analytic functions'' \cite[p.~206]{Jacobson1937}. His algebras
could be over any field. They did not have to be associative or finite
dimensional. So the definition could use only the product rule.

\begin{definition}[{\cite[p.~206]{Jacobson1937}}]\label{def:basic:derivation}
Let $A$ be an algebra over a field $k$. A \emph{derivation} of $A$ is a
$k$-linear map $D\colon A \to A$ with
\[
D(xy) = D(x)\,y + x\,D(y) \qquad (x, y \in A).
\]
\end{definition}

We list two changes that other authors made to this definition. Kharchenko
works with rings and writes maps on the right \cite[p.~7]{Kharchenko1991}.
Hochschild lets the values lie in a larger algebra $B$ that contains $A$ \cite[p.~688]{Hochschild1942}. A linear map
$D\colon A \to B$ with the same rule is then a derivation of $A$ into $B$
\cite[Definition~3.1, p.~688]{Hochschild1942}. We use the wider setting when
we say so.

We ask next what follows from the definition alone. Jacobson answers with the
Leibniz formula, which gives the $k$-th power of a derivation on a product
\cite[formula~(7), p.~207]{Jacobson1937}:
\[
D^k(xy) = \sum_{i=0}^{k} \binom{k}{i}\, D^{k-i}(x)\, D^{i}(y).
\]
For commutative rings, Freudenburg adds two more rules (Proposition~\ref{prop:basic:rules}).

\begin{proposition}[{\cite[Proposition~1.6(b),(c), p.~13]{Freudenburg2017}}]\label{prop:basic:rules}
Let $B$ be a commutative ring and $D \in \mathrm{Der}(B)$. For $t \in B$ and $n \ge 1$,
$D(t^n) = n\,t^{n-1}D(t)$. If $f, g \in B$ and $g$ is invertible, then
$D(f g^{-1}) = g^{-2}\bigl(g\,D(f) - f\,D(g)\bigr)$.
\end{proposition}

We call them the power rule and the quotient rule.

We turn to prime characteristic, where the formula takes a special form. Let the field have characteristic $p$. Then the middle binomial coefficients vanish, so
$D^p(xy) = D^p(x)\,y + x\,D^p(y)$. Jacobson states this as the fact that $D^p$
is a derivation \cite[formula~(8), p.~207]{Jacobson1937}. In characteristic $0$ the powers need not be derivations: for $D = d/dx$ on $k[x]$ we get $D^2(x^2) = 2$, while $D^2(x)\,x + x\,D^2(x) = 0$.

We now have a definition and a rule for repeated application.
Section~\ref{sec:basic-theory:der-as-lie-algebra} asks how the derivations of
one algebra fit together, and Jacobson answered it in his Theorem~1
\cite[p.~207]{Jacobson1937}.

\section{Physics: Poisson brackets and quantum derivations}\label{sec:basic-theory:physics}

The oldest application of the product rule is in physics. % claim-ok: own signpost connecting the moved section to this chapter
Dirac posed the
question in 1925, in the same paper in which he replaced the
Heisenberg matrix commutation by a Poisson bracket
\cite{Dirac1925}. He first writes the two laws that the classical Poisson
bracket obeys.

\begin{definition}[{\cite[equation after (10), p.~648]{Dirac1925}}]\label{def:basic:dirac-poisson}
Let $[x,y]$ be the Poisson bracket of classical mechanics. The two laws
IA and IIa of Dirac read
\[
[x+y,z]=[x,z]+[y,z],\qquad [xy,z]=x[y,z]+[y,z]x .
\]
\end{definition}

The second law is the product rule, so with $x$ fixed the map
$y\mapsto[x,y]$ is a derivation on the algebra of classical observables.
Dirac then declares that the same algebraic law holds for quantum
observables (Theorem~\ref{thm:basic:dirac-correspondence}) \cite[p.~648]{Dirac1925}.

\begin{theorem}[{\cite[equation (11), p.~648]{Dirac1925}}]\label{thm:basic:dirac-correspondence}
For any two quantum variables $x$ and $y$,
\[
xy-yx = \frac{ih}{2\pi}\,[x,y],
\]
where $h$ denotes the Planck constant and $[x,y]$ on the right is the Poisson
bracket of the corresponding classical quantities.
\end{theorem}

In particular the commutator $x\mapsto xy-yx$ is a scalar multiple of the
Poisson bracket, so it is again a derivation. With $H$ the Hamiltonian
function of the system, equation (14) of Dirac reads $\dot x=[x,H]$
\cite[equation (14), p.~651]{Dirac1925}: time evolution is an inner
derivation of the algebra of observables. The question of which derivations
of an operator algebra are inner is taken up in
Section~\ref{sec:banach-operator:kadison-sakai}. % claim-ok: own signpost to the chapter that now holds Kadison 1966

\hbfigure{02-basic-theory/dirac-correspondence}{The correspondence principle of Dirac. The classical bracket obeys laws IA and IIa \cite[equation after (10), p.~648]{Dirac1925}, so with $x$ fixed the map $y\mapsto[x,y]$ is a derivation. Equation (11) \cite[equation (11), p.~648]{Dirac1925} makes the quantum commutator a scalar multiple of that bracket, so it is a derivation too.}{fig:basic:dirac}

The Leibniz rule itself can be obtained rather than assumed. Hosseini derives
it in its most general form \cite[Abstract, p.~1]{Hosseini2022LeibnizRule}.
